
\subsubsection{Reproducibility Challenges in ML}

Recent empirical studies have quantified reproducibility gaps in ML research. Jiang et al. analyzed 348 bugs from PyTorch-based systems, finding that environment defects are often interface-related and many involve API mismatches. Wolter et al. surveyed ML reproducibility practices, reporting that the majority of repositories lack automated testing. Collberg and Proebsting found that only approximately half of published systems could be built from source, with environment configuration being a primary barrier.

While these studies characterize the problem space, they do not provide concrete interventions. Our work investigates whether library-level behavioral validation can detect certain classes of environment-stage defects.

\subsubsection{Dependency Management and Environment Isolation}

The standard approach to ML reproducibility relies on dependency pinning and environment isolation. However, version pinning addresses which versions are used, not whether those versions behave as expected. Two limitations motivate our approach: (1) pinning cannot validate behavioral invariants across adjacent versions, and (2) pinned environments do not detect assumption violations within a single version.

Our contracts complement version pinning by validating behavioral invariants that should hold regardless of specific versions.

\subsubsection{Integration Testing}

Repository-level integration tests provide regression detection for specific codebases. However, integration tests are repo-specific artifacts that encode usage patterns for particular projects rather than reusable library-level behavioral specifications. The key distinction: pytest validates whether code works in a specific repository, while contracts validate whether libraries behave as documented across repositories.

\subsubsection{Property-Based Testing and Metamorphic Testing}

Property-based testing frameworks (QuickCheck, Hypothesis) validate software by generating random inputs and checking specified properties. Metamorphic testing validates systems by asserting mathematical relations between inputs and outputs. Both approaches have been applied to ML systems.

We build on these foundations but differ in scope: (1) we target documented API invariants specifically in ML libraries, (2) we enforce execution constraints suitable for environment-stage validation, and (3) we deploy contracts at the library level for cross-repository reuse.

\subsubsection{Positioning}

Our approach addresses library-level behavioral validation tailored to ML API patterns, deployable at environment-setup time. We trade formal completeness for pragmatic deployability: contracts are lightweight Python decorators that validate documented invariants without requiring type-system integration.
